\section{The three-column \texorpdfstring{\(\mathrm{SU}(3)\)}{SU(3)}
Gram cone and its rank-loss \texorpdfstring{\(S^8\)}{S8}}
\label{sec:hr3c-main}

This section proves the topology forced by the left actions of
\(\mathrm U(3)\) and \(\mathrm{SU}(3)\) on three complex columns.  The
calculation produces a canonical topological \(S^8\), but it does not
equip that sphere with a complex structure, a fibration, a gauge bundle,
or a physical interpretation.

\subsection{Conventions}
\label{subsec:hr3c-conventions}

Put
\begin{equation}
 \mathscr V_3=\operatorname{Mat}_{3\times3}(\mathbb C),
 \qquad
 \lVert A\rVert_F^2=\operatorname{tr}(A^\dagger A).
 \label{eq:hr3c-source}
\end{equation}
This is a real vector space of dimension \(18\).  Both
\(\mathrm U(3)\) and \(\mathrm{SU}(3)\) act on the left:
\begin{equation}
 U\cdot A=UA.
 \label{eq:hr3c-left-action}
\end{equation}
Let
\begin{equation}
 \mathcal P_3=
 \{H\in\operatorname{Herm}_3(\mathbb C):H\succeq0\},
 \qquad
 \mathcal K_3=
 \{H\in\mathcal P_3:\operatorname{tr}H=1\}.
 \label{eq:hr3c-gram-cone-slice}
\end{equation}
Write \(S^{17}_F\subset\mathscr V_3\) for the Frobenius unit sphere and
\(\mathscr V_3^+=\mathscr V_3\cup\{\infty\}\cong S^{18}\) for the
one-point compactification, with \(\infty\) fixed by each group.

\subsection{The \texorpdfstring{\(\mathrm U(3)\)}{U(3)} orbit space}
\label{subsec:hr3c-u3}

\begin{theorem}[Exact Gram classification]
\label{thm:hr3c-u3-gram-classification}
The map
\begin{equation}
 \Gamma_U:\mathscr V_3\longrightarrow\mathcal P_3,
 \qquad
 \Gamma_U(A)=A^\dagger A
 \label{eq:hr3c-u3-gram-map}
\end{equation}
is \(\mathrm U(3)\)-invariant, is surjective, and separates
\(\mathrm U(3)\)-orbits.  It induces a homeomorphism
\begin{equation}
 \overline\Gamma_U:
 \mathscr V_3/\mathrm U(3)\xrightarrow{\ \sim\ }\mathcal P_3.
 \label{eq:hr3c-u3-quotient-homeomorphism}
\end{equation}
\end{theorem}

\begin{proof}
Invariance is the identity
\[
 (UA)^\dagger(UA)=A^\dagger U^\dagger UA=A^\dagger A.
\]
For \(H\succeq0\), its positive-semidefinite square root satisfies
\[
 (H^{1/2})^\dagger H^{1/2}=H,
\]
so the map is surjective.

Suppose \(A^\dagger A=B^\dagger B\).  Define
\[
 T:\operatorname{im}A\longrightarrow\operatorname{im}B,
 \qquad
 T(Ax)=Bx.
\]
If \(Ax=0\), then
\[
 \lVert Bx\rVert^2
 =x^\dagger B^\dagger Bx
 =x^\dagger A^\dagger Ax
 =\lVert Ax\rVert^2
 =0,
\]
so \(T\) is well-defined.  For \(x,y\in\mathbb C^3\),
\[
 \langle T(Ax),T(Ay)\rangle
 =x^\dagger B^\dagger By
 =x^\dagger A^\dagger Ay
 =\langle Ax,Ay\rangle.
\]
Thus \(T\) is a complex-linear isometry.  Extend it to a unitary
\(U\in\mathrm U(3)\).  Then \(UA=B\), so the Gram matrix separates
orbits.

Finally,
\[
 \lVert A\rVert_F^2=\operatorname{tr}(A^\dagger A).
\]
Hence \(\Gamma_U\) is proper.  The induced map
\(\overline\Gamma_U\) is a proper continuous bijection between locally
compact Hausdorff spaces and is therefore a homeomorphism.
\end{proof}

\begin{lemma}[Explicit ball chart on the trace-one Gram body]
\label{lem:hr3c-density-ball}
Let
\begin{equation}
 E=\{X\in\operatorname{Herm}_3(\mathbb C):\operatorname{tr}X=0\}
 \label{eq:hr3c-traceless-space}
\end{equation}
with Frobenius inner product.  Then
\(\dim_{\mathbb R}E=8\), and there is an explicit homeomorphism
\begin{equation}
 \chi:\mathcal K_3\xrightarrow{\ \sim\ }B(E)\cong B^8
 \label{eq:hr3c-chi-homeomorphism}
\end{equation}
which sends the singular locus in \(\mathcal K_3\) onto \(S(E)\).
\end{lemma}

\begin{proof}
Put \(c=I_3/3\).  If \(u\in S(E)\), then \(u\ne0\) and
\(\operatorname{tr}u=0\), so
\(\lambda_{\min}(u)<0\).  Define
\begin{equation}
 R(u)=-\frac{1}{3\lambda_{\min}(u)}>0.
 \label{eq:hr3c-radial-length}
\end{equation}
The eigenvalues of \(c+s u\) are
\(1/3+s\lambda_j(u)\).  Therefore, for \(s\geq0\),
\begin{equation}
 c+s u\succeq0
 \quad\Longleftrightarrow\quad
 0\leq s\leq R(u),
 \label{eq:hr3c-positive-radial-interval}
\end{equation}
and \(c+R(u)u\) is singular.

Every \(H\in\mathcal K_3\setminus\{c\}\) has the unique expression
\[
 H=c+s u,
 \qquad
 s=\lVert H-c\rVert_F,
 \qquad
 u=\frac{H-c}{\lVert H-c\rVert_F}.
\]
Define
\begin{equation}
 \chi(c)=0,
 \qquad
 \chi(c+s u)=\frac{s}{R(u)}u,
 \label{eq:hr3c-chi-formula}
\end{equation}
and
\begin{equation}
 \chi^{-1}(0)=c,
 \qquad
 \chi^{-1}(r u)=c+rR(u)u
 \quad(0<r\leq1).
 \label{eq:hr3c-chi-inverse}
\end{equation}
The minimum-eigenvalue function is continuous.  Since it is strictly
negative on the compact sphere \(S(E)\), \(R\) is continuous, finite,
positive, and bounded above and below.  Equations
\eqref{eq:hr3c-chi-formula} and \eqref{eq:hr3c-chi-inverse} are therefore
continuous also at the origin and are inverse homeomorphisms.
Equation \eqref{eq:hr3c-positive-radial-interval} shows that
\[
 \chi\bigl(\{H\in\mathcal K_3:\operatorname{rank}H\leq2\}\bigr)
 =S(E).
\]
\end{proof}

\begin{corollary}[Unit-sphere \(\mathrm U(3)\) quotient]
\label{cor:hr3c-u3-unit-quotient}
There is a homeomorphism of pairs
\begin{equation}
 \left(
 S^{17}_F/\mathrm U(3),
 \{[A]:\operatorname{rank}A\leq2\}
 \right)
 \cong
 (B^8,S^7).
 \label{eq:hr3c-u3-unit-pair}
\end{equation}
\end{corollary}

\begin{proof}
On \(S^{17}_F\),
\[
 \operatorname{tr}(A^\dagger A)=\lVert A\rVert_F^2=1.
\]
Apply Theorem~\ref{thm:hr3c-u3-gram-classification} and
Lemma~\ref{lem:hr3c-density-ball}.
\end{proof}

\begin{theorem}[Full \(\mathrm U(3)\) cone and compactification]
\label{thm:hr3c-u3-compactification}
There are homeomorphisms
\begin{equation}
 \mathscr V_3/\mathrm U(3)
 \cong\mathcal P_3
 \cong\mathbb R^8\times\mathbb R_{\geq0}
 \label{eq:hr3c-u3-half-space}
\end{equation}
and
\begin{equation}
 \mathscr V_3^+/\mathrm U(3)
 \cong
 \mathcal P_3^+
 \cong B^9.
 \label{eq:hr3c-u3-ball-nine}
\end{equation}
Under the second homeomorphism,
\begin{equation}
 \partial B^9
 =
 \{[A]:\operatorname{rank}A\leq2\}
 \cup\{[\infty]\}
 \cong S^8.
 \label{eq:hr3c-u3-boundary-s8}
\end{equation}
\end{theorem}

\begin{proof}
For \(H\ne0\), write
\[
 H=r\rho,
 \qquad
 r=\operatorname{tr}H>0,
 \qquad
 \rho=\frac{H}{\operatorname{tr}H}\in\mathcal K_3.
\]
Thus \(\mathcal P_3\) is the infinite cone on
\(\mathcal K_3\cong B^8\).

We give the compactification map explicitly.  Put
\[
 t=\frac{r-1}{r+1}\in(-1,1)
\]
and define
\begin{equation}
 \mathcal D^9=
 \{(y,t)\in E\oplus\mathbb R:
 \lVert y\rVert_F+|t|\leq1\}.
 \label{eq:hr3c-double-cone}
\end{equation}
Define
\begin{equation}
 \begin{split}
 \Xi_U(r\rho)
  &=\bigl((1-|t|)\chi(\rho),t\bigr),\\
 \Xi_U(0)&=(0,-1),\\
 \Xi_U(\infty)&=(0,1).
 \end{split}
 \label{eq:hr3c-xi-u}
\end{equation}
For \(|t|<1\), the inverse is
\begin{equation}
 x=\frac{y}{1-|t|},
 \qquad
 \rho=\chi^{-1}(x),
 \qquad
 r=\frac{1+t}{1-t},
 \qquad
 H=r\rho.
 \label{eq:hr3c-xi-u-inverse}
\end{equation}
As \(r\downarrow0\) or \(r\to\infty\), the factor \(1-|t|\) tends
uniformly to zero.  A subset of \(\mathcal P_3\) is unbounded if and only
if its trace is unbounded.  Thus \(\Xi_U\) is a homeomorphism
\(\mathcal P_3^+\to\mathcal D^9\).

For \(p=(y,t)\ne0\), put
\[
 \nu(p)=\sqrt{\lVert y\rVert_F^2+t^2},
 \qquad
 \gamma(p)=\lVert y\rVert_F+|t|
\]
and define
\begin{equation}
 \Theta(p)=\frac{\gamma(p)}{\nu(p)}p,
 \qquad
 \Theta(0)=0.
 \label{eq:hr3c-double-cone-ball-map}
\end{equation}
Then \(\lVert\Theta(p)\rVert_2=\gamma(p)\leq1\).  For nonzero
\(q\in B(E\oplus\mathbb R)\), the inverse is
\begin{equation}
 \Theta^{-1}(q)=
 \frac{\lVert q\rVert_2}
 {\lVert q_E\rVert_F+|q_{\mathbb R}|}\,q.
 \label{eq:hr3c-double-cone-ball-inverse}
\end{equation}
Hence \(\Theta:\mathcal D^9\to B^9\) is a homeomorphism and takes
\(\partial\mathcal D^9\) onto \(S^8\).

For finite \(H=r\rho\), Equation~\eqref{eq:hr3c-xi-u} lies on
\(\partial\mathcal D^9\) exactly when
\(\lVert\chi(\rho)\rVert_F=1\), equivalently when \(H\) is singular.
The points \(0\) and \(\infty\) are the two boundary tips.  This proves
\eqref{eq:hr3c-u3-boundary-s8}.

The canonical map
\[
 \mathscr V_3^+/\mathrm U(3)
 \longrightarrow
 (\mathscr V_3/\mathrm U(3))^+
\]
is a continuous bijection: the finite orbits agree and the added orbit is
\(\{\infty\}\).  Its source is compact and its target is Hausdorff, hence
it is a homeomorphism.  This proves
\eqref{eq:hr3c-u3-ball-nine}.

Finally, \(\Theta\circ\Xi_U\) sends \(\infty\) to
\(N=(0,\ldots,0,1)\in\partial B^9\).  The Cayley map
\begin{equation}
 (x',x_9)\longmapsto
 \left(
 \frac{2x'}{\lVert x'\rVert^2+(1-x_9)^2},
 \frac{1-\lVert x\rVert^2}
 {\lVert x'\rVert^2+(1-x_9)^2}
 \right)
 \label{eq:hr3c-cayley-map}
\end{equation}
is a homeomorphism
\[
 B^9\setminus\{N\}\longrightarrow
 \mathbb R^8\times\mathbb R_{\geq0}.
\]
Its inverse is
\begin{equation}
 (v,s)\longmapsto
 \left(
 \frac{2v}{\lVert v\rVert^2+(s+1)^2},
 \frac{\lVert v\rVert^2+s^2-1}
 {\lVert v\rVert^2+(s+1)^2}
 \right).
 \label{eq:hr3c-cayley-inverse}
\end{equation}
This proves \eqref{eq:hr3c-u3-half-space}.
\end{proof}

\subsection{The \texorpdfstring{\(\mathrm{SU}(3)\)}{SU(3)} orbit space}
\label{subsec:hr3c-su3}

\begin{theorem}[Gram matrix plus alternating determinant]
\label{thm:hr3c-su3-invariants}
Define
\begin{equation}
 \Gamma_{SU}:\mathscr V_3\longrightarrow
 \operatorname{Herm}_3(\mathbb C)\times\mathbb C,
 \qquad
 \Gamma_{SU}(A)=(A^\dagger A,\det A).
 \label{eq:hr3c-su-invariant-map}
\end{equation}
Its image is exactly
\begin{equation}
 \mathcal Y_3=
 \{(H,z)\in\mathcal P_3\times\mathbb C:
 |z|^2=\det H\}.
 \label{eq:hr3c-su-invariant-image}
\end{equation}
It separates \(\mathrm{SU}(3)\)-orbits and induces a homeomorphism
\begin{equation}
 \mathscr V_3/\mathrm{SU}(3)\xrightarrow{\ \sim\ }\mathcal Y_3.
 \label{eq:hr3c-su-quotient}
\end{equation}
\end{theorem}

\begin{proof}
The map is \(\mathrm{SU}(3)\)-invariant, and
\[
 \det(A^\dagger A)
 =\overline{\det A}\det A
 =|\det A|^2.
\]
Thus its image is contained in \(\mathcal Y_3\).

Conversely, take \((H,z)\in\mathcal Y_3\).  If \(H\) is singular, then
\(z=0\), and \(A=H^{1/2}\) realizes the pair.  If \(H\succ0\), put
\[
 \omega=\frac{z}{\sqrt{\det H}}\in S^1,
 \qquad
 Q_\omega=\operatorname{diag}(\omega,1,1)\in\mathrm U(3),
 \qquad
 A=Q_\omega H^{1/2}.
\]
Then \(A^\dagger A=H\) and \(\det A=z\).  This proves surjectivity.

Suppose
\[
 A^\dagger A=B^\dagger B,
 \qquad
 \det A=\det B.
\]
Theorem~\ref{thm:hr3c-u3-gram-classification} gives
\(U_0\in\mathrm U(3)\) with \(U_0A=B\).  If \(A\) is invertible, then
\[
 \det U_0=\frac{\det B}{\det A}=1,
\]
so \(U_0\in\mathrm{SU}(3)\).

If \(\operatorname{rank}A\leq2\), choose a unitary \(V\) which is the
identity on \(\operatorname{im}B\) and has determinant
\((\det U_0)^{-1}\) on \((\operatorname{im}B)^\perp\).  This is possible
because the orthogonal complement has positive complex dimension.  Then
\[
 VB=B,
 \qquad
 \det(VU_0)=1,
 \qquad
 (VU_0)A=B.
\]
Thus the invariant separates \(\mathrm{SU}(3)\)-orbits.

Since \(\operatorname{tr}H=\lVert A\rVert_F^2\), the invariant map is
proper.  The induced continuous bijection onto the closed Hausdorff set
\(\mathcal Y_3\) is therefore a homeomorphism.
\end{proof}

\begin{theorem}[The unit-sphere \(\mathrm{SU}(3)\) quotient]
\label{thm:hr3c-su-unit-sphere}
Let
\begin{equation}
 \mathcal Y_{3,1}=
 \{(H,z)\in\mathcal Y_3:\operatorname{tr}H=1\}.
 \label{eq:hr3c-su-unit-image}
\end{equation}
There is an explicit homeomorphism
\begin{equation}
 S^{17}_F/\mathrm{SU}(3)
 \cong
 \mathcal Y_{3,1}
 \xrightarrow{\ \sim\ }
 S(E\oplus\mathbb C)
 \cong S^9.
 \label{eq:hr3c-su-unit-sphere-result}
\end{equation}
\end{theorem}

\begin{proof}
The first homeomorphism follows from
Theorem~\ref{thm:hr3c-su3-invariants}.  To construct the second, let
\[
 x=\chi(H)\in B(E),
 \qquad
 r=\lVert x\rVert_F.
\]
If \(r<1\), then \(H\succ0\) and \(z\ne0\).  Define
\begin{equation}
 \Omega(H,z)=
 \left(
 x,\sqrt{1-r^2}\frac{z}{|z|}
 \right)
 \quad(r<1),
 \qquad
 \Omega(H,0)=(x,0)
 \quad(r=1).
 \label{eq:hr3c-omega}
\end{equation}
Its norm is exactly one:
\[
 \lVert x\rVert_F^2+
 \left|\sqrt{1-r^2}\frac{z}{|z|}\right|^2
 =r^2+(1-r^2)=1.
\]
The map is continuous on the full-rank locus.  As \(H\) approaches the
singular locus, \(r\to1\), so its complex coordinate tends to zero
regardless of phase.  Hence it is continuous everywhere.

Given \((x,w)\in S(E\oplus\mathbb C)\), set
\(H=\chi^{-1}(x)\).  If \(\lVert x\rVert_F<1\), then \(w\ne0\), and the
unique inverse value is
\begin{equation}
 z=\sqrt{\det H}\frac{w}{|w|}.
 \label{eq:hr3c-omega-inverse-z}
\end{equation}
If \(\lVert x\rVert_F=1\), then \(w=0\), \(H\) is singular, and \(z=0\).
Thus \(\Omega\) is a continuous bijection from a compact space to a
Hausdorff space and hence is a homeomorphism.
\end{proof}

\begin{theorem}[The weighted cone and its compactification]
\label{thm:hr3c-su-full-cone}
There is a homeomorphism
\begin{equation}
 \Phi_{SU}:\mathcal Y_3\xrightarrow{\ \sim\ }
 E\oplus\mathbb C\cong\mathbb R^{10}
 \label{eq:hr3c-su-r10}
\end{equation}
and therefore
\begin{equation}
 \mathscr V_3^+/\mathrm{SU}(3)\cong S^{10}.
 \label{eq:hr3c-su-s10}
\end{equation}
\end{theorem}

\begin{proof}
For \((H,z)\ne(0,0)\), set
\begin{equation}
 s=\sqrt{\operatorname{tr}H},
 \qquad
 h=\frac{H}{s^2},
 \qquad
 \zeta=\frac{z}{s^3}.
 \label{eq:hr3c-weighted-normalization}
\end{equation}
The exponents are forced by scalar multiplication in the original source:
for every real \(\lambda>0\),
\begin{equation}
 (\lambda A)^\dagger(\lambda A)=\lambda^2A^\dagger A,
 \qquad
 \det(\lambda A)=\lambda^3\det A,
 \qquad
 s(\lambda A)=\lambda s(A).
 \label{eq:hr3c-scaling-degrees}
\end{equation}
Moreover,
\[
 \operatorname{tr}h=1,
 \qquad
 |\zeta|^2
 =\frac{|z|^2}{s^6}
 =\frac{\det H}{s^6}
 =\det h.
\]
Thus \((h,\zeta)\in\mathcal Y_{3,1}\).  Define
\begin{equation}
 \Phi_{SU}(0,0)=0,
 \qquad
 \Phi_{SU}(H,z)=s\,\Omega(h,\zeta)
 \quad(s>0).
 \label{eq:hr3c-su-radial-homeomorphism}
\end{equation}
Since \(\Omega(h,\zeta)\) has norm one,
\[
 \lVert\Phi_{SU}(H,z)\rVert=s.
\]
For nonzero \(v\in E\oplus\mathbb C\), the inverse is obtained from
\[
 s=\lVert v\rVert,
 \qquad
 (h,\zeta)=
 \Omega^{-1}\left(\frac{v}{\lVert v\rVert}\right),
 \qquad
 H=s^2h,
 \qquad
 z=s^3\zeta.
\]
The norm identity proves continuity at the origin in both directions, so
\(\Phi_{SU}\) is a homeomorphism.

As in the proof of Theorem~\ref{thm:hr3c-u3-compactification}, compactness
of \(\mathrm{SU}(3)\) gives
\[
 \mathscr V_3^+/\mathrm{SU}(3)
 \cong
 (\mathscr V_3/\mathrm{SU}(3))^+.
\]
The one-point compactification of \(\mathbb R^{10}\) is \(S^{10}\).
\end{proof}

\subsection{Rank strata, the standard \texorpdfstring{\(S^8\)}{S8}, and
determinant winding}
\label{subsec:hr3c-rank-complement}

\begin{proposition}[Stabilizers and dimensions]
\label{prop:hr3c-rank-strata}
Let
\[
 \mathscr V_{3,(r)}
 =\{A\in\mathscr V_3:\operatorname{rank}_{\mathbb C}A=r\}.
\]
For \(A\in\mathscr V_{3,(r)}\),
\begin{equation}
 \operatorname{Stab}_{\mathrm U(3)}(A)\cong\mathrm U(3-r),
 \qquad
 \operatorname{Stab}_{\mathrm{SU}(3)}(A)\cong\mathrm{SU}(3-r),
 \label{eq:hr3c-stabilizers}
\end{equation}
with the rank-three and \(\mathrm{SU}(1)\) stabilizers understood to be
trivial.  The quotient-stratum dimensions are
\begin{center}
\begin{tabular}{c|c|c|c|c|c}
\(r\)&
\(\dim_{\mathbb R}\mathscr V_{3,(r)}\)&
\(\dim U(3)\!\cdot\!A\)&
\(\dim \mathscr V_{3,(r)}/U(3)\)&
\(\dim SU(3)\!\cdot\!A\)&
\(\dim \mathscr V_{3,(r)}/SU(3)\)\\ \hline
3&18&9&9&8&10\\
2&16&8&8&8&8\\
1&10&5&5&5&5\\
0&0&0&0&0&0
\end{tabular}
\end{center}
On the unit sphere, the nonzero quotient-stratum dimensions are
\[
 \begin{array}{c|ccc}
 r&3&2&1\\ \hline
 U(3)&8&7&4\\
 SU(3)&9&7&4.
 \end{array}
\]
The trace-one rank-one locus is \(\mathbb CP^2\).
\end{proposition}

\begin{proof}
A left stabilizer fixes \(\operatorname{im}A\) pointwise.  Its unrestricted
unitary action on \((\operatorname{im}A)^\perp\) gives
\(\mathrm U(3-r)\).  Imposing determinant one gives
\(\mathrm{SU}(3-r)\).  The real dimension of the complex rank-\(r\)
matrix stratum is
\[
 2r(6-r).
\]
Subtracting the orbit dimensions gives the table; intersecting a nonzero
rank stratum with the unit sphere subtracts one more dimension.

A trace-one positive-semidefinite matrix has rank one exactly when it has
the form \(vv^\dagger\) with \(\lVert v\rVert=1\).  The phase of \(v\)
does not change \(vv^\dagger\), so this locus is
\[
 S^5/S^1\cong\mathbb CP^2.
\]
\end{proof}

\begin{theorem}[The exact rank-loss pair and complement]
\label{thm:hr3c-rank-loss-pair}
Let
\begin{equation}
 \mathcal Y_{3,\mathrm{sing}}
 =\{(H,0)\in\mathcal Y_3:\operatorname{rank}H\leq2\}.
 \label{eq:hr3c-singular-invariant-locus}
\end{equation}
The homeomorphism \(\Phi_{SU}\) gives homeomorphisms of pairs
\begin{equation}
 (\mathcal Y_3,\mathcal Y_{3,\mathrm{sing}})
 \cong
 (E\oplus\mathbb C,E\oplus\{0\})
 \label{eq:hr3c-affine-pair}
\end{equation}
and
\begin{equation}
 (\mathcal Y_3^+,\mathcal Y_{3,\mathrm{sing}}^+)
 \cong
 (S^{10},S^8_{\mathrm{lin}}).
 \label{eq:hr3c-compact-pair}
\end{equation}
Consequently the compactified rank-loss locus is a standard,
topologically locally flat codimension-two \(S^8\subset S^{10}\), and
\begin{equation}
 S^{10}\setminus S^8
 \cong
 E\times(\mathbb C\setminus\{0\})
 \cong
 \mathbb R^8\times(\mathbb C\setminus\{0\})
 \simeq S^1.
 \label{eq:hr3c-complement}
\end{equation}
The map
\begin{equation}
 \vartheta(H,z)=\frac{z}{|z|}
 =\frac{\det A}{|\det A|}
 \label{eq:hr3c-phase-map}
\end{equation}
is the angular deformation-retraction coordinate and represents a
generator of
\[
 H^1(S^{10}\setminus S^8;\mathbb Z)
 \cong
 \pi_1(S^{10}\setminus S^8)
 \cong\mathbb Z
\]
after choosing the standard positive orientation on \(S^1\).
\end{theorem}

\begin{proof}
In Equation~\eqref{eq:hr3c-omega}, \(H\) is singular exactly when
\(\lVert\chi(H)\rVert_F=1\), in which case the complex coordinate is zero.
Conversely, a zero complex coordinate on the unit \(S^9\) has its
\(E\)-coordinate on \(S(E)\), and its inverse Gram matrix is singular.
Radial extension through Equation~\eqref{eq:hr3c-su-radial-homeomorphism}
therefore gives exactly \eqref{eq:hr3c-affine-pair}.

The homeomorphism \(\Phi_{SU}\) is proper and extends to one-point
compactifications, proving \eqref{eq:hr3c-compact-pair}.  The second member
is the one-point compactification of the linear pair
\[
 (E\oplus\mathbb C,E\oplus\{0\}).
\]
It is the standard topological pair \((S^{10},S^8_{\mathrm{lin}})\);
therefore the \(S^8\) is topologically locally flat at every point.

This local-flatness statement is not a smooth-stratification assertion.
The unit rank-deficient \(S^7\) contains the rank-one orbit stratum
\(\mathbb CP^2\), with rank-two complement
\(S^7\setminus\mathbb CP^2\).  The homeomorphism \(\chi\) uses the minimum
eigenvalue and need not be smooth across eigenvalue-collision strata.

Removing the linear subspace from the affine pair gives
\eqref{eq:hr3c-complement}.  Under \(\Phi_{SU}\), the phase of its complex
coordinate is \(z/|z|\), so \(\vartheta\) is the usual angular projection
\[
 E\times(\mathbb C\setminus\{0\})\longrightarrow S^1.
\]
This is a deformation retraction.  The explicit quotient loop
\begin{equation}
 A_\theta=
 \frac1{\sqrt3}\operatorname{diag}(e^{i\theta},1,1),
 \qquad
 0\leq\theta\leq2\pi,
 \label{eq:hr3c-generator-loop}
\end{equation}
satisfies
\[
 A_\theta^\dagger A_\theta=\frac13I_3,
 \qquad
 \det A_\theta=\frac{e^{i\theta}}{3\sqrt3}.
\]
Thus \(\vartheta([A_\theta])=e^{i\theta}\) winds once and is a generator.
\end{proof}

\begin{corollary}[Residual determinant-phase circle]
\label{cor:hr3c-residual-circle}
The residual group
\[
 \mathrm U(3)/\mathrm{SU}(3)\cong S^1
\]
acts on \(E\oplus\mathbb C\) by
\begin{equation}
 e^{i\theta}\cdot(x,w)=(x,e^{i\theta}w).
 \label{eq:hr3c-residual-action}
\end{equation}
Its fixed set after compactification is the standard
\(S^8\subset S^{10}\), and
\begin{equation}
 S^{10}/S^1\cong B^9,
 \qquad
 \operatorname{Fix}(S^1)=S^8=\partial B^9.
 \label{eq:hr3c-residual-quotient}
\end{equation}
\end{corollary}

\begin{proof}
If \(U\in\mathrm U(3)\), left multiplication fixes \(H=A^\dagger A\) and
sends \(z=\det A\) to \((\det U)z\).  Every phase occurs as a determinant.
Equation~\eqref{eq:hr3c-omega} therefore converts the residual action into
\eqref{eq:hr3c-residual-action}.  On the unit sphere,
\[
 S(E\oplus\mathbb C)/S^1
 \cong
 \{(x,|w|):\lVert x\rVert_F^2+|w|^2=1,\ |w|\geq0\}
 \cong B(E).
\]
On the affine cone,
\[
 (E\oplus\mathbb C)/S^1
 \cong E\times\mathbb R_{\geq0}.
\]
One-point compactification gives \(B^9\), and the fixed set is
\((E\oplus\{0\})^+\cong S^8\), which maps to the boundary.
\end{proof}

\subsection{The exact rank-two/rank-three distinction}
\label{subsec:hr3c-two-ladders}

\begin{theorem}[Quadratic Lorentz cone versus cubic rank-three cone]
\label{thm:hr3c-two-ladders}
For
\(\mathscr V_2=\operatorname{Mat}_{3\times2}(\mathbb C)\), the left
\(\mathrm{SU}(3)\)- and \(\mathrm U(3)\)-orbit relations coincide, and
\begin{equation}
 \mathscr V_2/\mathrm{SU}(3)
 \cong
 \operatorname{Herm}_2(\mathbb C)^{\succeq0}.
 \label{eq:hr3c-two-column-quotient}
\end{equation}
In the coordinates
\begin{equation}
 H_2(t,x_1,x_2,x_3)=
 \begin{pmatrix}
  t+x_3&x_1-i x_2\\
  x_1+i x_2&t-x_3
 \end{pmatrix},
 \label{eq:hr3c-lorentz-coordinates}
\end{equation}
\begin{equation}
 \operatorname{tr}H_2=2t,
 \qquad
 \det H_2=t^2-x_1^2-x_2^2-x_3^2,
 \label{eq:hr3c-lorentz-determinant}
\end{equation}
and
\begin{equation}
 H_2\succeq0
 \quad\Longleftrightarrow\quad
 t\geq0,\qquad
 t^2\geq x_1^2+x_2^2+x_3^2.
 \label{eq:hr3c-lorentz-positivity}
\end{equation}
Thus the two-column quotient is exactly the closed future
\(3+1\)-dimensional Lorentz cone; its one-point compactification is
\(B^4\), with rank-one-plus-infinity boundary \(S^3\).

By contrast, the three-column \(\mathrm U(3)\) quotient is the
nine-real-dimensional, Jordan-rank-three cone
\(\operatorname{Herm}_3(\mathbb C)^{\succeq0}\).  Its determinant is cubic,
its trace-one section is \(B^8\), and its compactified
rank-deficient-plus-infinity boundary is \(S^8\).  It is not a Lorentz
cone.
\end{theorem}

\begin{proof}
Every \(3\times2\) matrix has rank at most two, so the determinant-adjustment
step in the proof of Theorem~\ref{thm:hr3c-su3-invariants} converts every
left \(\mathrm U(3)\) equivalence into a left \(\mathrm{SU}(3)\)
equivalence.  The Gram proof of
Theorem~\ref{thm:hr3c-u3-gram-classification}, with two columns, gives
\eqref{eq:hr3c-two-column-quotient}.

Direct expansion of \eqref{eq:hr3c-lorentz-coordinates} gives
\eqref{eq:hr3c-lorentz-determinant}.  A \(2\times2\) Hermitian matrix is
positive semidefinite exactly when its trace and determinant are
nonnegative, which gives \eqref{eq:hr3c-lorentz-positivity}.  Its trace-one
section is
\[
 t=\frac12,
 \qquad
 x_1^2+x_2^2+x_3^2\leq\frac14,
\]
a closed \(3\)-ball.  The compactification argument of
Theorem~\ref{thm:hr3c-u3-compactification}, with \(E\) of dimension three,
gives \(B^4\) with boundary \(S^3\).

For three columns, Theorem~\ref{thm:hr3c-u3-compactification} gives the
\(B^9/S^8\) statements.  The determinant of a \(3\times3\) Hermitian
matrix is homogeneous of degree three, so this rank-three cone is not the
rank-two quadratic Lorentz cone.
\end{proof}

\subsection{The alternating invariant and the precise nonclaim}
\label{subsec:hr3c-nonclaim}

The extra invariant retained by the square-matrix
\(\mathrm{SU}(3)\) quotient is
\begin{equation}
 z=\det A
 =\varepsilon_{ijk}A_{i1}A_{j2}A_{k3}.
 \label{eq:hr3c-alternating-invariant}
\end{equation}
It is exactly the complex alternating volume of the three columns:
\[
 \det(UA)=\det(U)\det(A).
\]
It is invariant for \(U\in\mathrm{SU}(3)\), its phase is rotated by the
residual \(\mathrm U(3)/\mathrm{SU}(3)\), and it vanishes exactly when the
three columns are linearly dependent.

This invariant theory supplies exact carrier and orbit data for a physical
translation, but does not by itself complete that translation.  The next
construction must specify a quark field, colour-charge and baryon
observables, a principal \(\mathrm{SU}(3)\) bundle with connection and
characteristic classes, and an interacting Yang--Mills sector, together with
maps relating each object to Equation~\eqref{eq:hr3c-alternating-invariant}.
Likewise, Theorem~\ref{thm:hr3c-u3-compactification} supplies a distinguished
\(S^8\), while an almost-complex or integrable-complex structure, a
three-fibre structure, and analogues of the candidate \(S^6\) special-fibre
attaching maps remain exact construction problems rather than consequences
already proved by the compactification theorem.

The exact conclusion is the topological carrier theorem
\begin{equation}
 \boxed{
 \mathscr V_3^+/\mathrm U(3)\cong B^9,\quad
 \partial B^9\cong S^8,\quad
 \mathscr V_3^+/\mathrm{SU}(3)\cong S^{10},\quad
 S^{10}\setminus S^8\simeq S^1.
 }
 \label{eq:hr3c-final-carrier}
\end{equation}
Any gauge-theoretic or geometric use of this carrier requires additional
explicit descent data.
